% Merged draft: Henry's teacher-level sweet-spot paper + Saksham's Phase B
% supervision-mismatch diagnostic, combined per PAPER_MERGE_PLAN.md.
%
% Status: drafting only. Reconstructed from paper/SLFD_draft.pdf (no .tex source
% existed) + paper/phase_b_results_section.md + PHASE_B_PAPER_RESULT_CARD.md +
% results/positive_control/*. Numbers are transcribed as faithfully as possible
% from those sources; anywhere a source number was ambiguous or inconsistent
% across documents, it is flagged inline with a \todo{} note rather than
% silently resolved. Drop this into the target Overleaf/venue template; this
% file uses a plain article skeleton so it is easy to paste section-by-section.

\documentclass[11pt]{article}
\usepackage[margin=1in]{geometry}
\usepackage{booktabs}
\usepackage{graphicx}
\usepackage{amsmath}
\usepackage{hyperref}
\usepackage{xcolor}

\newcommand{\todo}[1]{\textcolor{red}{[TODO: #1]}}

\title{When Does Privileged Process Supervision Help?\\
A Teacher-Level Sweet Spot, a Verified Transfer Null,\\
and a Supervision-Mismatch Diagnosis}
\author{\todo{author list -- anonymize per venue if double-blind}}
\date{}

\begin{document}
\maketitle

\begin{abstract}
Process reward models (PRMs) score intermediate reasoning steps, and many
pipelines construct their step-level supervision automatically from a
stronger teacher model. We study two questions: when does privileged
(answer-aware) supervision help such a teacher, and separately, can that
advantage be distilled into a small, deployable, answer-blind student
verifier? On a controlled three-tier difficulty probe (GSM8K, MATH,
OlympiadBench), a privileged teacher's step-error localization improves only
in a tractability sweet spot: privilege is redundant on easy problems the
teacher can already self-verify (GSM8K, n.s.), unusable on the hardest
problems where the reference is too long to exploit (OlympiadBench, n.s.),
and helpful precisely in between (MATH, $+0.051$ F1, 95\% CI $[0.01, 0.09]$),
replicating across model families (Qwen-27B, $+0.082$ F1). This advantage does
not distill: a 1.5B student trained on the privileged teacher's labels is
statistically indistinguishable from one trained on no-GT labels, and neither
beats majority vote on downstream best-of-$N$ re-ranking. We show this null is
not explained by pipeline insensitivity -- a positive control confirms the
training path tracks teacher quality ($+0.050$ ROC-AUC, 95\% CI
$[0.021, 0.078]$, significant) -- nor primarily by student capacity: under an
identical training and evaluation path, a 3B verifier trained on
format-matched ProcessBench-style gold labels reaches $0.73$--$0.79$ ROC-AUC
across eight seeds on full ProcessBench-MATH (up to $0.83$ with post-hoc
score averaging across seeds), while the same generated teacher labels
remain weak ($0.55$--$0.63$ ROC-AUC) at both 1.5B and 3B scale. The evidence
points to a supervision distribution/format mismatch, not raw capacity, as the
dominant failure mode, and offers a practical diagnostic: test whether
format-matched supervision can support the target verifier before scaling
students or searching over nearby training recipes.
\end{abstract}

\section{Introduction}
\label{sec:intro}

Process reward models assign scores to intermediate reasoning steps and are
increasingly used to supervise or verify mathematical reasoning systems.
Because human step-level annotation is expensive, many PRM pipelines
construct supervision automatically, either by giving a labeling teacher
privileged (answer-aware) access to a reference solution, or by other
weak-supervision means. Two questions follow naturally, and most prior work
answers only the first: (1) does privileged access actually make the teacher
a better labeler, and if so, when; and (2) does that advantage survive
distillation into a small, deployable, answer-blind student verifier?

We answer both under a controlled setup that holds the student architecture,
optimizer, and evaluation benchmark fixed while varying exactly one thing at
a time. First, we characterize \emph{when} privileged supervision helps a
teacher (\S\ref{sec:sweetspot}): it is not uniformly useful, but pays off
only in a tractability window where the teacher both needs the reference (it
cannot already self-verify) and can act on it (the reference is short enough
to follow). Second, we ask whether this real, mechanistically-understood
teacher-level advantage distills into a small student PRM
(\S\ref{sec:transfer}): it does not, and we show this null is not an
artifact of an insensitive training pipeline. Third, having ruled out
pipeline insensitivity, we ask the natural follow-up -- is the failure to
distill a student-capacity problem or a supervision problem?
(\S\ref{sec:mismatch}). Under the identical training and evaluation path, we
show that format-matched gold process supervision \emph{does} produce a
competent verifier, while our generated teacher labels do not, at both 1.5B
and 3B student scale. This decomposes what would otherwise be a flat,
somewhat unsatisfying null into three connected findings: a real,
well-characterized teacher effect; a verified failure of that effect to
distill, with pipeline sensitivity confirmed; and a concrete diagnosis of
why it fails, pointing at supervision distribution/format rather than
capacity.

\section{Related Work}
\label{sec:related}

\subsection{Process supervision and process reward models}
A line of work has established that supervising \emph{how} a model reasons,
step by step, is more effective than supervising only whether the final
answer is correct. Lightman et al.~\cite{lightman2023verify} compared the two
regimes directly and found that process supervision substantially
outperforms outcome supervision on the challenging MATH dataset, releasing
the large PRM800K corpus of step-level human labels to support the finding.
Because human step annotation is expensive, Math-Shepherd~\cite{wang2023mathshepherd}
showed that process supervision can be constructed automatically -- assigning
a reward to each solution step without manual labels -- and used the
resulting signal both to re-rank candidate solutions and to drive step-level
reinforcement learning. These process reward models (PRMs) treat each step's
correctness as a scalar quantity and deploy the PRM as a frozen scorer or
ranker. Our work shares the step-level granularity but departs from the
scalar view for its teacher-level analysis: we treat the teacher's step
judgment, including its natural-language critique, as the object to be
distilled rather than a number to be predicted.

\subsection{Generative PRMs}
A more recent thread makes the verifier itself generative, producing a
reasoning trace before it judges. GenPRM~\cite{zhao2025genprm} performs
explicit chain-of-thought reasoning with code verification before scoring
each step, while ThinkPRM~\cite{khalifa2025thinkprm} fine-tunes a long-CoT
verifier that articulates a verification rationale for every step and
reaches strong ProcessBench performance from only a few thousand process
labels. These models are close to ours in spirit -- they generate critiques
rather than emit bare scores -- but they are designed as verifiers evaluated
at inference time, not as a teacher whose score-plus-critique behavior is
distilled into a separate, smaller, answer-blind student. None of them
isolates when privileged (answer-aware) judgment actually helps, which is
the question we take up in \S\ref{sec:sweetspot}.

\subsection{Weak and generated process supervision}
\label{sec:related-weak}
Recent weak-supervision approaches attempt to train PRMs without dense human
process labels, using outcome labels, generated rationales, or model-derived
signals~\cite{sun2025freeprm}. These methods make PRM training more
scalable, but they also raise the question we study in
\S\ref{sec:mismatch}: when generated labels fail, is the problem the
student, the recipe, or the supervision distribution? ProcessBench itself
frames earliest-error identification as a benchmark on which existing PRMs
often fail to generalize across sources~\cite{zheng2024processbench}.
Qwen2.5-Math PRMs are strong public baselines trained at much larger scale
and provide an external calibration point for our results
(\S\ref{sec:mismatch})~\cite{zhang2025qwenmathprm,qwen25mathprm800k}.
Concurrent work studies PRM transfer across GSM8K, MATH, OlympiadBench, and
OmniMath more directly -- ProcessLID reports explicit transfer tables across
these sources~\cite{bai2026processlid}, and RetrievalPRM frames PRM failures
as question/step out-of-distribution issues across the same
sources~\cite{zhu2025retrievalprm}. Our contribution relative to this thread
is not that cross-source ProcessBench transfer is itself novel, but the
controlled contrast between generated teacher labels and ProcessBench-style
gold labels under one fixed student, optimizer, and evaluation path
(\S\ref{sec:mismatch}).

\subsection{Expert--amateur contrast}
Our framing descends from work that extracts signal from the gap between a
stronger and a weaker model. CLEAR~\cite{rufail2025clear} contrasts feedback
from a larger ``expert'' and a smaller ``amateur'' model and applies the
refined difference to iteratively improve outputs across reasoning tasks; it
is the direct predecessor of the present project. LightReasoner~\cite{wang2025lightreasoner}
similarly exploits the behavioral divergence between a strong expert and a
weak amateur to pinpoint high-value reasoning moments and distill them,
notably without ground-truth labels, to sharpen reasoning. We retain the
ground-truth-free constraint at student test time but invert the axis of
interest: rather than contrasting models of different capacity, we study a
teacher with a privileged view -- access to the reference solution while
labeling -- distilled into a student that is answer-blind by construction.
The relevant axis for us is information, not size.

\subsection{Data-centric learning and privileged information}
Data-centric AI work argues that systematic changes in data quality can
dominate model-side changes~\cite{mazumder2022dataperf}. Our
\S\ref{sec:mismatch} result is a PRM-specific instance of that pattern: a
fixed small verifier's performance changes substantially with the
supervision source, holding the model and optimizer fixed. The
privileged-label experiments in \S\ref{sec:sweetspot}--\S\ref{sec:transfer}
are also related to learning using privileged
information~\cite{vapnik2009privileged}, but with an important boundary:
privileged information can improve the teacher's per-example inference
without producing labels that distill into a better answer-blind student.

\subsection{Our position}
We make three contributions, and we are deliberate about which is
load-bearing. First (\S\ref{sec:sweetspot}), we characterize \emph{when}
privileged supervision is worth anything at the teacher level: a
privilege~$\times$~difficulty~$\times$~richness effect in which a teacher's
answer-aware judgment improves step-error localization only when the problem
is hard enough that the teacher cannot self-verify \emph{and} the privilege
is rich enough to act on. To our knowledge this conditional structure --
privilege buys signal precisely where self-verification fails, and only when
the privileged information is a reasoning trace rather than a number -- has
not been isolated in the PRM or expert--amateur literature. This is the
validated result and the paper's spine. Second (\S\ref{sec:transfer}), we
show this advantage does not distill into a small student verifier, and we
rule out the most obvious confound -- that the training pipeline simply
cannot detect any teacher-quality difference -- with a positive control.
Third (\S\ref{sec:mismatch}), we diagnose the failure further: under an
identical training/evaluation path, format-matched gold process supervision
does produce a competent verifier while our generated labels do not, which
points at supervision distribution/format rather than student capacity as
the dominant explanation.

\section{Experimental Setup}
\label{sec:setup}

\subsection{Teacher-level privilege probe (\S\ref{sec:sweetspot})}
We evaluate a teacher's ability to localize the first erroneous step in a
multi-step solution under three supervision conditions that differ only in
the privileged information made available to the teacher at labeling time:
(i) \textbf{no-GT}, an answer-blind condition in which the teacher sees only
the problem and the candidate solution; (ii) \textbf{+answer}, in which the
teacher additionally sees the final reference answer; and (iii)
\textbf{+solution}, in which the teacher sees the full worked reference
solution. Performance is reported as ProcessBench-style first-error F1,
together with error-detection accuracy. Reported deltas are taken against the
no-GT condition, which serves as the unprivileged reference point.

\subsection{Student distillation (\S\ref{sec:transfer})}
The student is \texttt{Qwen2.5-1.5B-Instruct} with a lightweight
score-and-critique head, trained via LoRA on teacher-produced score and
natural-language critique labels ($N_\text{train}=1000$,
$N_\text{eval}=400$, 2 epochs). We compare students trained on privileged
(+solution) versus no-GT teacher labels, and additionally a
privileged-score-only ablation that drops the critique loss.

\subsection{Distribution-matched diagnostic (\S\ref{sec:mismatch})}
To separate student capacity from supervision distribution/format, we run a
second, matched training/evaluation path on \texttt{Qwen2.5-3B-Instruct}
with a score-only head, LoRA adapters, binary cross-entropy loss, balanced
batches, and 500 training steps. All headline results in
\S\ref{sec:mismatch} evaluate on the same shuffled ProcessBench-MATH split
of 1{,}000 solutions (6{,}505 total steps, 594 error steps), using exact
serial scoring (batch size 1). We compare two supervision sources under this
identical path: (a) ProcessBench-style gold step labels from the first 400
examples of a non-MATH source configuration (GSM8K or OmniMath), flattened
to per-step labels, and (b) our generated teacher labels from the
privileged/no-GT labeling pipeline described above. Exact problem-string
overlap between each source-config training file and the MATH-1000
evaluation split is zero.

\section{The Privilege Sweet Spot (Teacher Level)}
\label{sec:sweetspot}

\subsection{Teacher selection}
Before probing the effect of privilege, we fix the teacher by an
upper-bound bake-off in which every candidate is granted ground-truth access
(the with-GT ceiling, $N=30$). Table~\ref{tab:teacher-bakeoff} shows the
official Gemma-class checkpoint attains the strongest first-error F1
($0.909$) with perfect format compliance (zero parse failures) and the
highest throughput among the candidates; a comparably-sized Qwen variant
follows closely (F1 $0.873$). A reasoning-tuned candidate collapses under
the structured-output protocol, failing to parse on the large majority of
instances (parse-failure rate $0.846$, F1 $0.000$) -- a reminder that verbal
reasoning fluency does not by itself confer reliable structured step
adjudication. We adopt the Gemma checkpoint as the primary teacher and
retain the Qwen variant for the cross-family replication reported below.

\begin{table}[t]
\centering
\small
\begin{tabular}{lrrrr}
\toprule
Teacher & $F_1$ & correct acc & error acc & parse fail \\
\midrule
Gemma-4-26B (primary) & \textbf{0.909} & 1.000 & 0.833 & 0.000 \\
Qwen-27B & 0.873 & 0.917 & 0.833 & 0.000 \\
Qwen-35B-A3B & 0.839 & 1.000 & 0.722 & 0.019 \\
Qwen-4B-Thinking & 0.000 & 1.000 & 0.000 & 0.846 \\
\bottomrule
\end{tabular}
\caption{Teacher bake-off, with-GT ceiling ($N=30$).}
\label{tab:teacher-bakeoff}
\end{table}

\subsection{The sweet spot}
Our central finding is that privileged supervision is \emph{conditionally}
beneficial: it helps only when the problem is hard enough that the teacher
cannot self-verify and tractable enough that the teacher can still exploit
the reference. The solution-privilege gap is non-significant on easy
problems (GSM8K, $-0.052$, n.s.), opens to a significant $+0.051$ F1 on
hard-but-tractable problems (MATH, 95\% CI $[0.01, 0.09]$), and returns to
non-significant on the hardest tier (OlympiadBench, $-0.03$, n.s.). The
three points trace an inverted U in difficulty, but we do not lean on the
two non-significant endpoints alone to assert a curve; the curvature is
corroborated by the mechanism analysis in \S\ref{sec:mechanism}, where the
rescue rate declines monotonically with reference length. The reading the
data support is conjunctive: privilege pays off only where the teacher both
needs the reference and can use it.

Table~\ref{tab:sweetspot} reports the three tiers in full, together with the
Qwen-27B cross-family replication. On GSM8K, neither the bare answer nor the
full solution improves an already-saturated judgment; the teacher is, in
effect, its own oracle, and external privilege is redundant. On MATH, where
unaided self-verification degrades, the full worked solution lifts
first-error F1 from $0.726$ to $0.777$ ($+0.051$, 95\% CI $[0.01, 0.09]$,
significant) and error recall from $0.608$ to $0.661$. A smaller preliminary
MATH run ($N=50$) points the same way with a larger nominal gap ($\sim
+0.13$ F1). OlympiadBench, harder still, returns the gap to $\approx 0$: the
references are too long for the teacher to follow, so the privilege it is
handed is no longer usable. The same inverted U faintly reappears within
MATH by difficulty level (strongest near the middle, weakest at the easiest
and hardest levels), corroborating that tractability is the moderator; but
the load-bearing evidence is the three-tier contrast, not the noisier
per-level split. Critically, the cross-family Qwen-27B teacher reproduces
the effect on MATH ($+0.082$ F1, error detection $0.526 \to 0.667$), which
substantially weakens any single-model-artifact explanation.

\begin{table}[t]
\centering
\small
\begin{tabular}{llrrrl}
\toprule
Tier & Benchmark (teacher) & no-GT & +answer & +solution & $\Delta$ (solution) \\
\midrule
Easy & GSM8K (Gemma, $N=50$) & 0.889 & 0.856 & 0.837 & $-0.052$ (n.s.) \\
Hard & MATH (Gemma, $N=400$) & 0.726 & 0.726 & 0.777 & $\mathbf{+0.051}$ $[0.01, 0.09]$ \\
Hardest & OlympiadBench (Gemma, OE) & -- & -- & -- & $-0.03$ (n.s.) \\
\midrule
Hard & MATH (Qwen-27B, $N=150$) & 0.653 & 0.690 & 0.735 & $+0.082$ \\
\bottomrule
\end{tabular}
\caption{The privilege sweet spot across difficulty tiers, with cross-family
replication. Values are first-error F1. MATH (Gemma) is the $N=400$
definitive run (95\% CI shown); OlympiadBench is OE-only, verified.}
\label{tab:sweetspot}
\end{table}

\subsection{Richness, not mere access}
The MATH panel isolates a second, sharper effect. The +answer condition --
a correct final answer but no reasoning -- is indistinguishable from the
answer-blind baseline: the two rows are identical, and the bare answer flips
zero predictions (measured at $N=150$; we verified the labeling pipeline is
correctly wired, so the null effect is behavioral, not an artifact). Only
the full worked solution moves the metric. The operative variable is
therefore not whether the teacher is privileged, but whether the privilege
is \emph{rich}: a reference reasoning trace supplies the step-by-step
scaffold needed to localize an error, whereas a final answer, which the
teacher cannot back-propagate into the faulty step, supplies nothing it can
act on.

\subsection{Mechanism: privilege rescues self-verification failures}
\label{sec:mechanism}
The sweet spot is not a raw difficulty effect; it is a rescue of
self-verification failures, decomposing into two gates ($N=400$). \textbf{Gate
1 (the teacher must need help).} Of 227 problems containing an error, the
no-GT teacher misses 89. Adding the worked solution rescues 29 of those
($33\%$), while breaking only 17 of the 138 it had already caught ($12\%$) --
the benefit concentrates exactly where self-verification fails. \textbf{Gate
2 (the teacher must be able to use it).} The rescue rate falls monotonically
with the length of the reference solution (short $0.37 \to$ mid $0.33 \to$
long $0.28$), which is the mechanism behind the collapse at the hardest
tier. Together the gates identify the moderator as
self-verification-failure~$\times$~tractability, precisely the inverted U
observed both across and within datasets.

\begin{table}[t]
\centering
\small
\begin{tabular}{p{2.4cm}p{4.6cm}p{5.4cm}}
\toprule
Problem (MATH level) & What no-GT missed & +solution feedback at the flagged step \\
\midrule
Tens digit of $7!+8!+\cdots+2006!$ (L3) & Predicted no error anywhere (missed step 3) & ``Incorrectly assumes the tens digit of a sum is simply the sum of the tens digits; fails to account for a carry from the units place.'' \\
Regular-tetrahedron centroid ratio (L5) & Flagged step 1 instead of the actual error at step 4 & ``Incorrectly concludes $PQ/AQ=1/3$ from a 3{:}1 median-division claim; in a tetrahedron the centroid divides the altitude so that $AP/PQ=3{:}1$, not $PQ/AQ$.'' \\
Base-6 representation of 355 (L2) & Predicted no error anywhere (missed step 0) & ``The division $355 \div 6$ is incorrect; $355 = 59 \times 6 + 1$, so the remainder should be 1.'' \\
\bottomrule
\end{tabular}
\caption{Three rescue cases from \texttt{results/evidence\_pack\_n400/per\_sample.jsonl}
(gold error step known, no-GT prediction wrong, +solution prediction correct),
selected to span error types: a missed carry/place-value step, a misapplied
geometric ratio, and a missed arithmetic slip. 37 such rescue cases exist in
the $N{=}400$ evidence pack; these three are illustrative, not
cherry-picked for effect size.}
\label{tab:flip-cases}
\end{table}

\section{Does the Teacher Advantage Transfer to the Student?}
\label{sec:transfer}

The teacher-level sweet spot in \S\ref{sec:sweetspot} is the validated
result. The natural next question is whether it distills: if we train a
small (1.5B) student PRM on the privileged teacher's score-plus-critique
labels, does that student inherit a usable advantage over one trained on
no-GT labels? Across a validated $N=1000$ run (1{,}000 training problems,
400 evaluation problems, Gemma-4 labeling confirmed through $\sim$32k served
requests), the answer is no. We report this as a clean, diagnosed negative,
not a dead end.

\subsection{Step-level transfer}
We compare three students -- privileged+critique, privileged score-only, and
no-GT+critique -- on threshold-free ROC-AUC, the appropriate measure here
because first-error F1 at a fixed score-head cutoff moves with calibration
rather than capability. Table~\ref{tab:step-transfer} shows the no-GT
student is highest on ROC-AUC ($0.641$ vs.\ $0.631$), and the run is
non-degenerate (predicted-error rate $\approx 0.3$ across all cells, so no
silent collapse inflates the comparison). Privilege does not produce a
better student verifier.

\begin{table}[t]
\centering
\small
\begin{tabular}{lrrrr}
\toprule
Student (teacher labels) & ROC-AUC & PR-AUC & $F_1$ & error recall \\
\midrule
Privileged + critique & 0.631 & 0.130 & 0.170 & 0.357 \\
Privileged, score-only & 0.624 & 0.135 & 0.184 & 0.427 \\
No-GT + critique & \textbf{0.641} & 0.121 & 0.198 & 0.467 \\
\bottomrule
\end{tabular}
\caption{Step-level transfer ($N_\text{train}=1000$, $N_\text{eval}=400$).
The no-GT student is highest on ROC-AUC -- privilege does not transfer.}
\label{tab:step-transfer}
\end{table}

\subsection{Downstream best-of-$N$}
The motivating promise of a ground-truth-free step-error localizer is a
test-time gain: given $N$ sampled solutions, the student PRM scores each and
we re-rank, selecting the highest-scoring candidate. Table~\ref{tab:bon}
reports the $N=1000$ result. The no-GT verifier re-ranks \emph{better} than
the privileged one ($0.373$ vs.\ $0.349$), and decisively neither beats the
majority-vote baseline ($\approx 0.39$). The oracle ($0.517$) shows
substantial headroom remains, so the failure is the verifiers', not the
candidate pool's.

\begin{table}[t]
\centering
\small
\begin{tabular}{lrrrr}
\toprule
Verifier & pass@1 & prm rerank & majority vote & oracle pass@$N$ \\
\midrule
No-GT student & 0.337 & \textbf{0.373} & 0.391 & 0.517 \\
Privileged student & 0.335 & 0.349 & 0.382 & 0.518 \\
\bottomrule
\end{tabular}
\caption{Downstream best-of-$N$ re-ranking ($N=1000$ problems, 8 candidates,
\texttt{gemma-2-9b-it} generator). The no-GT verifier re-ranks better;
neither beats majority vote.}
\label{tab:bon}
\end{table}

\subsection{Why the null: a diffuse, not absent, teacher gap}
\label{sec:diffuse}
The negative is not that privileged and no-GT labels are identical. A probe
of the actual Gemma-4 labeling teacher reproduces a real $+0.07$
solution gap (F1 $0.716 \to 0.786$; error recall $0.577 \to 0.654$, $N=150$
MATH) -- consistent with \S\ref{sec:sweetspot}. The issue is the
\emph{shape} of the gap. Over 6{,}340 labeled steps, privileged and no-GT
labels agree on only $69\%$; on the $\sim 31\%$ where they differ, the
disagreement is nearly symmetric: privilege flags an error where no-GT does
not on 1{,}001 steps, while no-GT flags one where privilege does not on
956. The teacher's privilege churns labels in both directions; its small net
quality gain is diffuse, not a clean directional signal a 1.5B student can
latch onto. A same-pool paired best-of-$N$ confirms the consequence: the
privileged and no-GT re-rankers are statistically indistinguishable
(McNemar $p=0.14$; 5 vs.\ 12 discordant pairs), both below majority vote.
This $N=200$ test is underpowered (17 discordant pairs). A shared-pool
$N=1000$ rerun was attempted locally on 2026-06-24: it retrained both
students and completed the ProcessBench re-eval step, but the paired
best-of-$N$ generation step itself failed on an HTTP 401 from the local
generation server (\texttt{localhost:8000}), which was serving the Gemma-4
labeling teacher under \texttt{launchd} rather than the lightweight
generator the script expected, and for which the background job's
environment did not carry a valid \texttt{OMLX\_API\_KEY}. This is a
configuration issue, not a result -- the run needs to be pointed at a
correctly-authenticated generation endpoint and re-launched. Note for
whoever does this: that attempt's freshly-retrained checkpoints scored
lower (ROC-AUC $0.55$/$0.58$) than the checkpoints behind
Tables~\ref{tab:step-transfer}--\ref{tab:bon} ($0.63$/$0.64$), most likely
because it used a shorter step budget (\texttt{MAX\_STEPS=2000}) than the
original training run; retrain with matching hyperparameters before running
the paired test, or the new $N=1000$ number will not be comparable to the
rest of this section.

\subsection{Positive control: the pipeline is sensitive to teacher quality}
\label{sec:positive-control}
The two results above leave one obvious confound open: does the training
pipeline fail to distill privilege specifically, or is it simply
insensitive to \emph{any} teacher-quality difference? We test this directly.
We relabel the same candidate solutions with a deliberately weaker labeling
model (\texttt{Qwen2.5-3B-Instruct}, in place of the 26B Gemma-4 teacher)
and train an identically-configured 1.5B student on the resulting labels.
This weak-teacher student reaches ROC-AUC $0.575$ on the full 1{,}000-problem
ProcessBench-MATH split (6{,}505 steps), compared to $0.625$ for the
Gemma-4-labeled privileged+critique student evaluated on the identical
6{,}505-step split.\footnote{This $0.625$ differs slightly from the $0.631$
reported for the same condition in Table~\ref{tab:step-transfer}: Table~\ref{tab:step-transfer}'s
$0.631$ is a 400-problem-eval snapshot committed 2026-06-18, whereas the
checkpoint/result path \texttt{results/ablation/priv\_critique/} was
subsequently retrained and re-evaluated at least twice more (an interim
re-run logged $0.629$ on 2026-06-22; a full-1000-eval re-run committed
2026-06-23 is the $0.625$ used here), each time overwriting the same file
path without a distinct version tag. An internal log entry from
2026-06-22 cites a $0.624$ comparison point for this same positive control,
which does not exactly match any committed \texttt{priv\_critique} snapshot
we can find, and is closest to the \emph{privileged score-only} student's
$0.624$ (Table~\ref{tab:step-transfer}) -- likely a labeling slip in that
log rather than a different run. We do not attempt to reconcile this
retroactively; instead we recompute the comparison directly from the
currently-committed, identically-sized (6{,}505-step) per-step score files
for both conditions so the number in this section is independently
reproducible from the repository as it stands.} A paired bootstrap
(10{,}000 resamples, same procedure as \S\ref{sec:mismatch-main}) on this
ROC-AUC gap, recomputed from the currently-committed per-step scores, gives
a mean difference of $+0.0497$, 95\% CI $[0.0206, 0.0783]$, two-sided
$p=0.0012$ -- significant, excluding zero, and consistent in direction and
magnitude with the originally logged estimate ($+0.048$,
$[0.020, 0.077]$). The weak-teacher student's score head is also markedly
less calibrated (it predicts ``error'' on only $1.6\%$ of steps, versus
$\approx 30\%$ for the Gemma-4-labeled cells), so part of this gap reflects
outright degeneracy rather than a smooth quality gradient; we report the
comparison as-is rather than adjusting for it.

This is the paper's key validity result for \S\ref{sec:transfer}: because
student quality tracks teacher quality when the teacher-quality difference
is large and unprivileged (weak vs.\ strong labeling model), the training
pipeline is demonstrably capable of transmitting a teacher-quality signal
into the student. The null in \S\ref{sec:diffuse} is therefore not explained
by pipeline insensitivity. It is specifically the privileged-vs-no-GT
distinction -- a much subtler, more diffuse signal than weak-vs-strong
teacher quality -- that fails to distill at this student scale and training
recipe. This motivates the next question directly: is the remaining gap a
matter of student capacity, or of the generated-label supervision itself?
We take this up in \S\ref{sec:mismatch}.

\section{Diagnosing the Null: Capacity or Supervision?}
\label{sec:mismatch}

Section~\ref{sec:transfer} leaves a specific question open. Having ruled out
pipeline insensitivity (\S\ref{sec:positive-control}), is the failure of
privilege to distill because the 1.5B student lacks the capacity to
represent the signal, or because the generated-label supervision itself is
poorly matched to what the ProcessBench-MATH evaluation needs? We address
this with a controlled ablation on \texttt{Qwen2.5-3B-Instruct}
(\S\ref{sec:setup}): holding the verifier architecture, training path, and
evaluation benchmark fixed, we vary only the source and format of step-level
supervision.

\subsection{Format-matched gold labels transfer; generated labels remain weak}
\label{sec:mismatch-main}
Table~\ref{tab:mismatch-main} shows the main result. With the student,
training path, and evaluation held fixed, ProcessBench-style labels from
GSM8K and OmniMath transfer strongly to full ProcessBench-MATH. Generated
teacher labels are much weaker, even when the generated-label baseline is
chosen as the strongest checkpoint from a broader preliminary sweep
(\S\ref{sec:mismatch-ablations}). Every GSM8K and OmniMath seed
significantly beats both generated-label baselines under sequence-cluster
bootstrap (resampling whole solutions rather than individual steps, which
are correlated); against the best generated-label baseline overall, gaps
range from $+0.093$ to $+0.155$ ROC-AUC across the eight seeds, all
two-sided $p = 0.0010$ (Table~\ref{tab:mismatch-gaps}).

\begin{table}[t]
\centering
\small
\resizebox{\textwidth}{!}{%
\begin{tabular}{llrrrrr}
\toprule
Training source & Seeds & ROC-AUC & PR-AUC & Best $F_1^*$ & Fixed $F_1$ & Pred error rate \\
\midrule
GSM8K gold $\to$ MATH1000 & 0--3 & 0.7515 (0.7256--0.7760) & 0.2188 (0.1935--0.2317) & 0.3144 & 0.2503 & 0.5111 \\
OmniMath gold $\to$ MATH1000 & 0--3 & 0.7694 (0.7539--0.7869) & 0.2524 (0.2277--0.2877) & 0.3276 & 0.3055 & 0.2920 \\
Qwen2.5-Math-7B-PRM800K (public) & -- & 0.8379 & 0.3254 & 0.3991 & 0.3953 & 0.1436 \\
Generated privileged labels, BCE & 0 & 0.5503 & 0.0992 & 0.1803 & 0.0000 & 0.0000 \\
Generated no-GT labels, rank loss & 0 & 0.6324 & 0.1418 & 0.2221 & 0.2151 & 0.3819 \\
\bottomrule
\end{tabular}%
}
\caption{Full ProcessBench-MATH ($N=1000$) transfer results under an
identical Qwen2.5-3B score-head training/eval path. Parentheses show seed
ranges. $^*$Best $F_1$ is threshold-swept on the eval slice; treat as
diagnostic (Table~\ref{tab:mismatch-calibration} reports held-out
calibrated $F_1$). Qwen2.5-Math-7B-PRM800K is an external calibration
baseline, not a matched training recipe -- our result is not a SOTA PRM
claim.}
\label{tab:mismatch-main}
\end{table}

\begin{table}[t]
\centering
\small
\resizebox{\textwidth}{!}{%
\begin{tabular}{llrl}
\toprule
Model A & Model B & ROC-AUC gap & 95\% CI ($p$) \\
\midrule
GSM8K gold seed 0 & best generated-label baseline & $+0.111$ & $[0.087, 0.134]$ ($0.0010$) \\
OmniMath gold seed 0 & best generated-label baseline & $+0.121$ & $[0.096, 0.146]$ ($0.0010$) \\
GSM8K gold seed 0 & generated privileged BCE & $+0.193$ & $[0.164, 0.221]$ ($0.0010$) \\
OmniMath gold seed 0 & generated privileged BCE & $+0.203$ & $[0.176, 0.232]$ ($0.0010$) \\
Qwen PRM800K & OmniMath gold seed 3 (best single seed) & $+0.051$ & $[0.033, 0.068]$ ($0.0010$) \\
\bottomrule
\end{tabular}%
}
\caption{Representative sequence-cluster bootstrap ROC-AUC gaps (2{,}000
resamples). Full 16-row table across all seeds is in the appendix
(\texttt{paper/generated/phase\_b\_*.tex}).}
\label{tab:mismatch-gaps}
\end{table}

Held-out threshold calibration (threshold chosen on the first 200 MATH
sequences, evaluated on the remaining 800) agrees with the ranking result:
GSM8K and OmniMath gold seeds reach calibrated $F_1$ of $0.27$--$0.33$,
against $0.17$ for generated privileged BCE and $0.21$ for the best
generated-label baseline (Table~\ref{tab:mismatch-calibration}).

\begin{table}[t]
\centering
\small
\begin{tabular}{lrrrr}
\toprule
Model & Calibrated $F_1$ & Eval ROC-AUC & Eval PR-AUC & Pred error rate \\
\midrule
GSM8K gold seed 0 & 0.3166 & 0.7452 & 0.2339 & 0.2488 \\
OmniMath gold seed 0 & 0.3062 & 0.7483 & 0.2292 & 0.2742 \\
Qwen2.5-Math-7B-PRM800K & 0.3855 & 0.8372 & 0.3242 & 0.1091 \\
Generated privileged BCE & 0.1673 & 0.5490 & 0.1002 & 0.3212 \\
Best generated-label baseline & 0.2085 & 0.6314 & 0.1431 & 0.1695 \\
\bottomrule
\end{tabular}
\caption{Held-out threshold calibration (threshold on first 200 MATH
sequences, evaluated on remaining 800).}
\label{tab:mismatch-calibration}
\end{table}

\subsection{The source distribution still matters}
This is not a blanket claim that any non-MATH process labels transfer
cleanly to MATH. OlympiadBench is high-variance in our runs: two seeds give
$0.585$ and $0.716$ ROC-AUC on full MATH-1000 (mean $0.651$), one
significantly below the best generated-label baseline
($-0.047$ ROC-AUC, $p=0.0010$) and the other significantly above it
($+0.084$, $p=0.0010$). We treat OlympiadBench as a boundary diagnostic
rather than a third clean headline source; the pattern still supports a
source-distribution view, but the source split matters.

\subsection{The failure is not explained by the score-head machinery alone}
\label{sec:mismatch-probes}
A plausible concern is that the 3B score-head architecture is simply too
weak to learn step-level mathematical verification. Table~\ref{tab:mismatch-probes}
argues against this as the primary explanation. Training the same score-head
path on the first 200 shuffled MATH examples and evaluating on the next 200
reaches $0.7129$ ROC-AUC -- not a benchmark claim, since train and eval come
from the same benchmark family, but evidence that the machinery can learn
the target when labels and distribution are aligned. Frozen-representation
probes tell the same story: class-balanced logistic probes on frozen
Qwen2.5-3B step-boundary embeddings trained on generated teacher labels
transfer weakly ($0.52$--$0.56$ ROC-AUC), while the matched
ProcessBench-gold probe reaches $0.631$. The failure is not merely that the
student representation cannot expose the target; the supervision
distribution is doing meaningful work.

\begin{table}[t]
\centering
\small
\begin{tabular}{llr}
\toprule
Diagnostic & Condition & ROC-AUC \\
\midrule
Score head, matched ProcessBench-MATH 200/200 & Gold labels & 0.7129 \\
Frozen probe, generated teacher labels & Privileged & 0.5497 \\
Frozen probe, generated teacher labels & No-GT & 0.5559 \\
Frozen probe, ProcessBench-gold 200/200 & Gold labels & 0.6308 \\
\bottomrule
\end{tabular}
\caption{Diagnostics separating score-head/representation capacity from
supervision mismatch.}
\label{tab:mismatch-probes}
\end{table}

\subsection{Training-recipe changes did not rescue generated labels}
\label{sec:mismatch-ablations}
Before reframing the result as a supervision mismatch, we tested a range of
score-head and distillation variants on the generated labels
(Table~\ref{tab:mismatch-ablations}). Naive score-plus-critique distillation
was unstable until a feedback-token truncation bug was fixed; after the fix,
equal-weight critique distillation remained weak or negative for privileged
labels. Pure ranking loss and BCE-plus-rank did not rescue privileged
transfer -- they often favored the no-ground-truth labels instead. Longer BCE
training increased variance without turning the generated privileged labels
into a competent MATH verifier. These negative ablations do not prove no
generated-label method can work, but combined with the 1.5B ablations in
\S\ref{sec:transfer} (which tried soft-BCE, hard-verdict, and logit-KD
distillation with the same result), they reduce the likelihood that the
observed gap is a single bad optimizer setting at a single student scale.

\begin{table}[t]
\centering
\small
\resizebox{\textwidth}{!}{%
\begin{tabular}{lrrl}
\toprule
Generated-label recipe & Priv ROC-AUC & No-GT ROC-AUC & Note \\
\midrule
Score+critique MSE, early run & 0.4582 & 0.4870 & collapsed near all-negative \\
Score-only MSE & 0.4844 & 0.4776 & stable but weak \\
Feedback-token truncation fixed & 0.4514 & 0.5093 & stable but negative for priv. \\
BCE, error weight 3 & 0.5411 & 0.4700 & best short privileged ablation, still weak \\
Rank-only, balanced batches & 0.4939 & 0.6306 & no-GT wins strongly \\
BCE+rank, balanced batches & 0.5214 & 0.6140 & no-GT wins strongly \\
BCE, error weight 3, longer training & 0.4782 & 0.4971 & longer BCE degrades/nulls signal \\
\bottomrule
\end{tabular}%
}
\caption{Preliminary generated-label ablations on the Qwen2.5-3B score-head
path (condensed; full 10-row table in the appendix).}
\label{tab:mismatch-ablations}
\end{table}

\subsection{Interpretation}
Taken together with \S\ref{sec:transfer}, the picture is: the training
pipeline is sensitive to teacher quality (\S\ref{sec:positive-control}); the
score-head machinery and student representation can learn a competent
verifier when supervision is distribution-matched
(\S\ref{sec:mismatch-main}--\ref{sec:mismatch-probes}); and a broad sweep of
loss functions and student scales fails to rescue the generated-label
signal specifically (\S\ref{sec:transfer}, \S\ref{sec:mismatch-ablations}).
The most parsimonious explanation across all three tracks of evidence is a
supervision distribution/format mismatch in the generated-label pipeline,
not a capacity ceiling in the student.

\section{Limitations}
\label{sec:limitations}

\textbf{Confounded comparison in \S\ref{sec:mismatch}.} The gold-vs-generated
contrast varies provenance (human-curated ProcessBench annotations vs.\
LLM-generated labels), source-problem distribution (GSM8K/OmniMath vs.\
MATH-style traces), and label format simultaneously. This paper should not
be read as isolating human-curated versus generated provenance. The missing
controlled cell is generated teacher labels on the \emph{same}
GSM8K/OmniMath source problems used by the gold-label training rows --
running this is the highest-value next experiment if time allows before
submission.

\textbf{Seed and scale asymmetry.} The generated-label baselines in
\S\ref{sec:mismatch} use one seed each from a $\sim$10-configuration
preliminary sweep, while the gold-source rows use four seeds; additional
generated-label seeds could narrow the gap, though the breadth of the
preliminary sweep (Table~\ref{tab:mismatch-ablations}) makes a pure
unlucky-seed explanation less likely. Separately, the 1.5B (\S\ref{sec:transfer})
and 3B (\S\ref{sec:mismatch}) tracks are not a controlled capacity ablation
of one another -- they differ in loss recipe and evaluation $N$ -- so the
observation that generated-label transfer stays weak at both scales is
suggestive, not a clean single-variable ablation of student capacity.

\textbf{Not a SOTA PRM claim.} The public Qwen2.5-Math-7B-PRM800K baseline
is stronger than any single gold-source seed on the same split
(Table~\ref{tab:mismatch-main}). The contribution is diagnostic: a
format-matched cheap check can reveal supervision mismatch before investing
in a larger student or a broader recipe search, not a new state-of-the-art
verifier.

\textbf{OlympiadBench and the positive control both carry caveats.}
OlympiadBench is high-variance over two seeds and is treated as a boundary
diagnostic, not a headline source. The positive control in
\S\ref{sec:positive-control} uses a weak-teacher student whose score head is
also miscalibrated (near-zero error-prediction rate), so part of the
measured gap may reflect outright degeneracy rather than a smooth
teacher-quality gradient.

\textbf{Downstream validation of \S\ref{sec:mismatch} is incomplete.} Unlike
\S\ref{sec:transfer}, the 3B gold-vs-generated result is evaluated only on
ProcessBench step-level metrics; we do not show the gold-supervised 3B
verifier beats majority vote or the generated-label student on downstream
best-of-$N$ re-ranking.

\section{Conclusion}
\label{sec:conclusion}

Under a fixed 3-tier teacher probe, privileged (answer-aware) process
supervision helps a teacher precisely in a tractability window -- where it
is needed and usable -- and this effect replicates across model families.
That advantage does not distill into a small student verifier: a 1.5B
student trained on privileged labels is statistically indistinguishable from
one trained on no-GT labels, and neither clears majority vote on downstream
best-of-$N$. This null is not explained by an insensitive training pipeline,
which we confirm with a positive control. Diagnosing further, we show that
under an identical training/evaluation path, format-matched gold process
supervision produces a competent verifier while our generated teacher labels
do not, at both 1.5B and 3B student scale -- pointing at supervision
distribution/format mismatch, not student capacity, as the dominant
explanation. For efficient PRM development, this suggests a practical
ordering: before scaling students or re-running the same distillation loop,
test whether the supervision distribution itself can support the target
verifier.

\begin{thebibliography}{99}

\bibitem{lightman2023verify} Lightman, H., Kosaraju, V., Burda, Y., Edwards,
H., Baker, B., Lee, T., Leike, J., Schulman, J., Sutskever, I., and Cobbe, K.
(2023). Let's Verify Step by Step. \emph{arXiv:2305.20050}.

\bibitem{wang2023mathshepherd} Wang, P., Li, L., Shao, Z., Xu, R., Dai, D.,
Li, Y., Chen, D., Wu, Y., and Sui, Z. (2023). Math-Shepherd: Verify and
Reinforce LLMs Step-by-step without Human Annotations. \emph{arXiv:2312.08935}.

\bibitem{zhao2025genprm} Zhao, J., et al.\ (2025). GenPRM: Scaling
Test-Time Compute of Process Reward Models via Generative Reasoning.
\emph{arXiv:2504.00891}.

\bibitem{khalifa2025thinkprm} Khalifa, M., et al.\ (2025). Process Reward
Models That Think. \emph{arXiv:2504.16828}.

\bibitem{rufail2025clear} Rufail, A., Kim, D., O'Brien, S., and Zhu, K.
(2025). CLEAR: Contrasting Textual Feedback with Experts and Amateurs for
Reasoning. \emph{arXiv:2504.07116}.

\bibitem{wang2025lightreasoner} Wang, J., Chen, Y., Li, Z., and Huang, C.
(2025). LightReasoner: Can Small Language Models Teach Large Language
Models Reasoning? \emph{arXiv:2510.07962}.

\bibitem{zheng2024processbench} Zheng, C., Zhang, Z., Zhang, B., Lin, R.,
Lu, K., Yu, B., Liu, D., Zhou, J., and Lin, J. (2024). ProcessBench:
Identifying Process Errors in Mathematical Reasoning. \emph{arXiv:2412.06559}.

\bibitem{zhang2025qwenmathprm} Zhang, Z., Zheng, C., Wu, Y., Zhang, B.,
Lin, R., Yu, B., Liu, D., Zhou, J., and Lin, J. (2025). The Lessons of
Developing Process Reward Models in Mathematical Reasoning.
\emph{arXiv:2501.07301}.

\bibitem{qwen25mathprm800k} Qwen Team. Qwen2.5-Math-7B-PRM800K model card.
\url{https://huggingface.co/Qwen/Qwen2.5-Math-7B-PRM800K}.

\bibitem{sun2025freeprm} Sun, L., Liu, C., Ma, X., Yang, T., Lu, W., and
Wu, N. (2025). FreePRM: Training Process Reward Models without Ground
Truth Process Labels. \emph{arXiv:2506.03570}.

\bibitem{bai2026processlid} Bai, L., Yin, F., Sun, J., and Chang, K.-W.
(2026). ProcessLID: Step-wise Internal Reward in LLM Reasoning.
OpenReview 5O5AlNVAbs.

\bibitem{zhu2025retrievalprm} Zhu, J., Zheng, C., Lin, J., Du, K., Wen, Y.,
Yu, Y., Wang, J., and Zhang, W. (2025). Retrieval-Augmented Process Reward
Model for Generalizable Mathematical Reasoning. \emph{arXiv:2502.14361}.

\bibitem{mazumder2022dataperf} Mazumder, M., et al.\ (2022). DataPerf:
Benchmarks for Data-Centric AI Development. \emph{arXiv:2207.10062}.

\bibitem{vapnik2009privileged} Vapnik, V. and Vashist, A. (2009). A New
Learning Paradigm: Learning Using Privileged Information. \emph{Neural
Networks}, 22(5--6):544--557.

\end{thebibliography}

\end{document}
